\section*{Bab \ref{ch:b1:atmospheric-circulation} --- Sirkulasi umum atmosfer}
% values: calc/ch28.py

\begin{solution}{exo:b1:atmospheric-circulation:1}
(a)~$f = 2\Omega\sin\varphi$: $0$ di khatulistiwa, \qty{7.29e-5}{s^{-1}} pada \ang{30},
\qty{1.031e-4}{s^{-1}} pada \ang{45} dan \qty{1.458e-4}{s^{-1}} di kutub.
(b)~$fV = \num{1.031e-4} \times 10 = \qty{1.03e-3}{m.s^{-2}}$, sekitar \num{1.05e-4}
dari $g$. (c)~Di Belahan Utara defleksinya ke kanan arah gerak: angin yang berhembus ke utara
didorong ke timur.
\end{solution}

\begin{solution}{exo:b1:atmospheric-circulation:2}
(a)~$|\nabla p|/\rho = (400/\num{4e5})/1.225 = \qty{8.16e-4}{m.s^{-2}}$.
(b)~$10/\num{8.16e-4} = \qty{12250}{s}$, \qty{3.4}{h}. (c)~Percepatan Coriolis
angin \qty{10}{m/s} pada \ang{45} LU adalah \qty{1.03e-3}{m.s^{-2}}: keduanya berukuran
sama (rasio $0.79$). Dalam beberapa jam setelah mulai bergerak, udara merasakan gaya Coriolis yang mampu
mengimbangi gradien tekanan: angin menetap dekat kelajuan geostrofik, di sini
$\num{8.16e-4}/\num{1.031e-4} = \qty{7.9}{m/s}$, alih-alih
dipercepat menuruni gradien.
\end{solution}

\begin{solution}{exo:b1:atmospheric-circulation:3}
(a)~$f = 2 \times \num{7.292e-5} \times \sin\ang{40} = \qty{9.374e-5}{s^{-1}}$;
$|\vec v_g| = (400/\num{2.5e5})/(1.225 \times \num{9.374e-5}) = \qty{13.9}{m/s}$, atau
$13.93/0.5144 = \qty{27.1}{kn}$. (b)~Di Belahan Utara angin menjaga tekanan rendah di sebelah kirinya: dengan tekanan rendah di selatan, ia berhembus ke barat, dari timur
(\omterm{def:b1:spherical-astronomy:horizontal}{azimut} \ang{90}). (c)~Tangkai menunjuk ke timur dari lingkaran stasiun, ke arah angin;
\qty{27.1}{kn} dibulatkan menjadi \qty{25}{kn}, dua bulu penuh dan setengah bulu, digambar di sisi searah jarum jam tangkai, di sini sisi selatannya, ke arah rendah.
\end{solution}

\begin{solution}{exo:b1:atmospheric-circulation:4}
(a)~\ang{15} LU: \omterm{def:b1:atmospheric-circulation:winds}{angin pasat} timur laut; \ang{15} LS: \omterm{def:b1:atmospheric-circulation:winds}{angin pasat} tenggara;
\ang{45} LU dan \ang{45} LS: \omterm{def:b1:atmospheric-circulation:winds}{angin barat}. (b)~Pada \ang{15} udara permukaan \omterm{def:b1:atmospheric-circulation:cells}{sel Hadley} kembali
menuju khatulistiwa dan dibelokkan ke barat oleh gaya Coriolis (ke kanan di utara, ke kiri di selatan): timur, $u$ negatif.
Pada \ang{45} udara permukaan \omterm{def:b1:atmospheric-circulation:cells}{sel Ferrel} bergerak ke arah kutub dan dibelokkan ke timur: barat,
$u$ positif. (c)~Dari ketiga \omterm{def:b1:spherical-astronomy:geographic}{lintang} itu, tekanan tertinggi berada di \ang{30} LU,
\qty{1016.8}{hPa}: \omterm{def:b1:atmospheric-circulation:cells}{tekanan tinggi subtropik} berada dekat \ang{30} LU (profil lengkap menempatkan
maksimumnya dekat \ang{35} LU).
\end{solution}

\begin{solution}{exo:b1:atmospheric-circulation:5}
(a)~$|\nabla p|/\rho = (200/\num{5e5})/1.225 = \qty{3.27e-4}{m.s^{-2}}$, dibagi dengan
$f$: \qty{2.6}{m/s} pada \ang{60}, \qty{2.9}{m/s} pada \ang{50}, \qty{4.5}{m/s} pada
\ang{30}, \qty{6.5}{m/s} pada \ang{20}, \qty{12.9}{m/s} pada \ang{10} dan
\qty{25.7}{m/s} pada \ang{5}. (b)~Gaya Coriolis memutar angin dalam waktu berorde
$1/f$: \qty{2.2}{h} pada \ang{60}, \qty{11}{h} pada \ang{10}, \qty{22}{h} pada \ang{5}, dan
tak terhingga lama di khatulistiwa. Dekat khatulistiwa medan tekanan berubah, dan gesekan bekerja,
jauh sebelum angin dapat dibelokkan sepanjang isobar: keseimbangan itu tidak pernah terbentuk, dan udara mengalir sebagian menuruni gradien, dengan kelajuan sedang.
\end{solution}

\begin{solution}{exo:b1:atmospheric-circulation:6}
(a)~$u = 464.6\,\sin^2\varphi/\cos\varphi$: \qty{14.2}{m/s} pada \ang{10},
\qty{32.2}{m/s} pada \ang{15}, \qty{57.8}{m/s} pada \ang{20}, \qty{134.1}{m/s} pada
\ang{30}. (b)~$u = 464.6\,(\cos^2\ang{10} - \cos^2\varphi)/\cos\varphi$:
\qty{17.7}{m/s} pada \ang{15}, \qty{42.9}{m/s} pada \ang{20}, \qty{117.9}{m/s} pada
\ang{30}. (c)~Udara yang mulai lebih jauh dari khatulistiwa mulai lebih dekat ke sumbu,
dengan momentum sudut yang lebih sedikit untuk dibagi: ia tiba dengan \omterm{def:b1:atmospheric-circulation:winds}{angin barat} yang lebih lemah pada setiap
\omterm{def:b1:spherical-astronomy:geographic}{lintang}. Menyebarkan cabang naik menurunkan \omterm{def:b1:atmospheric-circulation:jet}{arus jet} yang diprediksi, meskipun tidak cukup
untuk mencapai \qty{40}{m/s} yang teramati.
\end{solution}

\begin{solution}{exo:b1:atmospheric-circulation:7}
(a)~Angin berhembus ke timur laut. Dengan membelakanginya, menghadap timur laut,
rendah berada di sebelah kiri: di barat laut. (b)~Angin berhembus ke
tenggara; di Belahan Selatan rendah berada di sebelah kanan orang yang membelakangi angin: di barat daya.
(c)~Angin permukaan diputar dari isobar menuju rendah, yaitu berlawanan arah jarum jam dari \omterm{def:b1:atmospheric-circulation:geostrophic}{angin geostrofik} di Belahan Utara.
\omterm{def:b1:atmospheric-circulation:geostrophic}{Angin geostrofik}, dan rendah yang tegak lurus di sebelah kirinya, karena itu diputar searah jarum jam dari jawaban (a): rendah
berada di antara barat laut dan utara, pada \omterm{def:b1:spherical-astronomy:horizontal}{azimut} $\ang{315} + \alpha$ untuk sudut melintas $\alpha$.
\end{solution}

\begin{solution}{exo:b1:atmospheric-circulation:8}
(a)~Luas antara khatulistiwa dan \ang{30} LU adalah $2\pi a_\oplus^2\sin\ang{30} =
\qty{1.275e14}{m^2}$; massa di atas setiap meter persegi antara kedua aras adalah
$\Delta p/g = \num{9e4}/9.81 = \qty{9174}{kg}$; totalnya adalah
\qty{1.170e18}{kg}. (b)~$\num{1.170e18}/\num{8.7e10} = \qty{1.34e7}{s}$, sekitar
\num{156} hari, lima bulan. (c)~$\num{1.170e18}/\num{21.4e10} = \qty{5.47e6}{s}$,
sekitar \num{63} hari, dua bulan.
\end{solution}

\begin{solution}{exo:b1:atmospheric-circulation:9}
(a)~$\num{0.01} \times \num{86400} = \qty{864}{m}$ sehari. (b)~$0.864 \times 9.8 =
\qty{8.5}{K}$ sehari. (c)~Udara mendingin dengan memancarkan radiasi inframerah: atmosfer
kehilangan lebih banyak energi melalui radiasi daripada yang diserapnya, dan turunnya udara, yang menghangatkannya karena kompresi, mengimbangi kerugian itu. (d)~Udara yang turun menghangat dan mengering
relatif terhadap kejenuhan, awannya menguap dan hujan tidak dapat terbentuk: sabuk \omterm{def:b1:atmospheric-circulation:cells}{tekanan tinggi subtropik} adalah sabuk langit cerah dan sedikit hujan, tempat gurun panas terbentuk di darat.
\end{solution}

\begin{solution}{exo:b1:atmospheric-circulation:10}
(a)~$V^2 + frV - fr|\vec v_g| = 0$ dengan $fr = \qty{55.85}{m/s}$:
$V = \tfrac12\bigl(-55.85 + \sqrt{55.85^2 + 4 \times 55.85 \times 15}\bigr) =
\qty{12.3}{m/s}$, \qty{82}{\percent} dari \omterm{def:b1:atmospheric-circulation:geostrophic}{angin geostrofik}. (b)~$V^2 - frV +
fr|\vec v_g| = 0$ hanya memiliki akar riil bila $(fr)^2 \geq 4fr|\vec v_g|$, yaitu
$|\vec v_g| \leq fr/4 = \qty{14.0}{m/s}$. (c)~Dekat pusat tekanan tinggi, tempat $r$ kecil,
batas $fr/4$ kecil: gradien tekanan harus lemah, isobar berjauhan dan angin ringan. Depresi tidak memiliki batas seperti itu: isobarnya dapat berdesakan dekat pusat, dengan angin kuat.
\end{solution}

\begin{solution}{exo:b1:atmospheric-circulation:11}
(a)~Untuk $\varphi$ kecil, $\sin\varphi \approx \varphi$ dan $\cos\varphi \approx 1$,
sehingga $u \approx \Omega a_\oplus\varphi^2$. (b)~$\ang{10} = \qty{0.1745}{rad}$:
$464.6 \times 0.1745^2 = \qty{14.15}{m/s}$, dibanding \qty{14.23}{m/s} secara eksak,
\qty{0.5}{\percent} lebih kecil. (c)~Dengan menyelesaikan $464.6\sin^2\varphi/\cos\varphi = 27.0$
secara numerik diperoleh $\varphi = \ang{13.7}$. (d)~Cincin itu sudah cocok dengan \omterm{def:b1:atmospheric-circulation:jet}{arus jet}
rerata tahunan yang sungguhan pada kurang dari \ang{14}: kekekalan momentum sudut menggambarkan cabang atas \omterm{def:b1:atmospheric-circulation:cells}{sel Hadley} hanya di tropik dalam; ke arah kutub dari sekitar
\ang{15}, pusaran mengambil alih.
\end{solution}

\begin{solution}{exo:b1:atmospheric-circulation:12}
(a)~Rerata $31.7/12 = \ang{2.6}$ LU; rentang dari \ang{12.8} LS hingga \ang{10.9} LU, yaitu
\ang{23.7}. (b)~\omterm{def:b1:spherical-astronomy:geographic}{Lintang} paling utara terjadi pada September, tiga bulan setelah
solstis Juni; yang paling selatan pada Februari, dua bulan setelah solstis Desember.
Jeda itu berasal dari penyimpanan panas, terutama di \omterm{def:b1:fluid-envelopes-structure:ocean-layers}{lapisan campuran} samudra
(\cref{ch:b1:fluid-envelopes-structure}), yang menghangat dan mendingin perlahan. (c)~Rerata zonal
merata-ratakan semua \omterm{def:b1:spherical-astronomy:geographic}{bujur}: jika di atas benua pita hujan berayun jauh
ke belahan musim panas, rerata semua \omterm{def:b1:spherical-astronomy:geographic}{bujur} dapat melintasi khatulistiwa meskipun pita di atas satu samudra tidak.
\end{solution}

\begin{solution}{pb:b1:atmospheric-circulation:1}
\begin{enumerate}
  \item $\rho = p/(R_dT) = \num{100600}/(287.05 \times 283.15) =
        \qty{1.238}{kg.m^{-3}}$.
  \item $|\nabla p| = \qty{400}{Pa}/\qty{2.5e5}{m} = \qty{1.6e-3}{Pa.m^{-1}}$.
  \item $|\nabla p|/\rho = \num{1.6e-3}/1.238 = \qty{1.293e-3}{m.s^{-2}}$.
  \item $\Omega = 2\pi/\num{86164.2} = \qty{7.292e-5}{rad.s^{-1}}$;
        $f = 2\Omega\sin\ang{50} = \qty{1.117e-4}{s^{-1}}$.
  \item $|\vec v_g| = \num{1.293e-3}/\num{1.117e-4} = \qty{11.57}{m/s}$.
  \item $11.57/0.514444 = \qty{22.5}{kn}$; $11.57 \times 3.6 = \qty{41.7}{km/h}$.
  \item Waktunya adalah $|\vec v_g|/(|\nabla p|/\rho)$ dan $|\vec v_g| = (|\nabla
        p|/\rho)/f$, sehingga ia adalah $1/f = \qty{8951}{s}$, \qty{2.49}{h}; dan $2\pi/f =
        \qty{15.6}{h}$. Pola yang bertahan sehari atau lebih karena itu mendekati keseimbangan.
  \item \omterm{def:b1:atmospheric-circulation:pgf}{Gaya gradien tekanan} menunjuk ke utara, ke arah rendah. \omterm{def:b1:atmospheric-circulation:geostrophic}{Angin geostrofik} berhembus sepanjang isobar dengan rendah di sebelah kirinya: ke timur,
        dari barat, \omterm{def:b1:spherical-astronomy:horizontal}{azimut} \ang{270}: \omterm{def:b1:atmospheric-circulation:winds}{angin barat}.
  \item Dengan membelakangi \omterm{def:b1:atmospheric-circulation:winds}{angin barat}, menghadap ke timur, sebelah kiri adalah utara, tempat rendah berada: hukum itu terpenuhi.
  \item Tangkai menunjuk ke barat dari lingkaran stasiun, ke arah datangnya angin; \qty{22.5}{kn} dibulatkan menjadi \qty{20}{kn}, dua bulu penuh; mereka berdiri di sisi searah jarum jam tangkai, sisi utaranya, ke arah rendah.
  \item Pada \ang{50} LS, $f < 0$: angin menjaga rendah di sebelah kanannya, sehingga ia berhembus ke barat, dari timur.
  \item $f(\ang{30}) = \qty{7.292e-5}{s^{-1}}$: \qty{17.7}{m/s}; $f(\ang{60}) =
        \qty{1.263e-4}{s^{-1}}$: \qty{10.2}{m/s}.
  \item $f(\ang{5}) = \qty{1.271e-5}{s^{-1}}$ memberi \qty{101.7}{m/s}, angin yang tidak pernah teramati dengan gradien seperti itu. Waktu pembelokan $1/f$ adalah \qty{22}{h}:
        keseimbangan tidak punya waktu untuk terbentuk, dan rumusnya tidak berlaku.
  \item $9.0/11.57 = 0.778$; sepanjang isobar $9.0\cos\ang{25} = \qty{8.16}{m/s}$,
        melintasi mereka menuju rendah $9.0\sin\ang{25} = \qty{3.80}{m/s}$.
  \item \omterm{def:b1:atmospheric-circulation:geostrophic}{Angin geostrofik} berhembus menuju \omterm{def:b1:spherical-astronomy:horizontal}{azimut} \ang{90}; diputar \ang{25} menuju rendah
        di utara, angin permukaan berhembus menuju \ang{65}, dari \ang{245}, barat-barat daya.
  \item $\tan\alpha = 1/(f\tau)$ memberi $\tau = 1/(\num{1.117e-4}\tan\ang{25}) =
        \qty{1.92e4}{s}$, \qty{5.3}{h}; model memprediksi $11.57\cos\ang{25} =
        \qty{10.5}{m/s}$, lebih besar daripada \qty{9.0}{m/s} yang diamati: gesekan sungguhan
        tidak tepat linear terhadap angin, dan memperlambatnya lebih dari model pada sudut ini.
  \item Melalui silinder berjari-jari $R = \qty{500}{km}$ dan tinggi $h =
        \qty{1}{km}$: $\dot m = 2\pi R h\rho u_r = 2\pi \times \num{5e5} \times 1000
        \times 1.238 \times 3.80 = \qty{1.48e10}{kg/s}$.
  \item $w = \dot m/(\rho\pi R^2) = 2hu_r/R = 2 \times 1000 \times 3.80/\num{5e5} =
        \qty{0.0152}{m/s}$, sekitar \qty{1310}{m} sehari.
  \item $1.31 \times 9.8 = \qty{12.9}{K}$ sehari: udara segera mencapai kejenuhan,
        dan depresi terisi awan dan hujan. Di bawah \omterm{def:b1:atmospheric-circulation:cyclone}{antisiklon} udara permukaan mengalir
        ke luar, udara di atasnya turun dan menghangat, dan langit menjadi cerah.
  \item $\Omega a_\oplus = \num{7.292e-5} \times \num{6.371e6} = \qty{464.6}{m/s}$
        di khatulistiwa; $464.6\cos\ang{30} = \qty{402.3}{m/s}$ pada \ang{30} LU.
  \item $u = 464.6\,\sin^2\varphi/\cos\varphi$: \qty{57.8}{m/s} pada \ang{20} LU dan
        \qty{134.1}{m/s} pada \ang{30} LU.
  \item $u = 464.6\,(\cos^2\ang{10} - \cos^2\ang{30})/\cos\ang{30} =
        \qty{117.9}{m/s}$.
  \item $134.1/42.3 = 3.17$. Pusaran \omterm{def:b1:spherical-astronomy:geographic}{lintang} menengah, depresi dan \omterm{def:b1:atmospheric-circulation:cyclone}{antisiklon},
        mengambil momentum sudut dari cabang atas dan membawanya ke tempat lain; dan udara
        sel itu naik di atas sabuk \omterm{def:b1:spherical-astronomy:geographic}{lintang} alih-alih diam di khatulistiwa. Rerata zonal juga
        merata-ratakan \omterm{def:b1:atmospheric-circulation:jet}{arus jet} yang berliku atas semua \omterm{def:b1:spherical-astronomy:geographic}{bujur}.
  \item $M/M_0 = (402.3 + 42.3)\cos\ang{30}/464.6 = 0.829$: udara telah kehilangan
        \qty{17}{\percent} dari momentum sudut yang akan ia pertahankan sebagai cincin bebas.
  \item \omterm{def:b1:atmospheric-circulation:geostrophic}{Angin geostrofik} di P adalah \qty{11.6}{m/s}, dari barat; angin yang diukur
        di permukaan adalah \qty{9.0}{m/s}, sekitar empat perlima darinya, dan diputar
        \ang{25} menuju rendah oleh gesekan.
\end{enumerate}
\end{solution}
